\chapter{E1: Softmax, MLP và CNN}
\section{Dataset và tiền xử lý}
Fashion-MNIST gồm 60k ảnh train và 10k test, 10 lớp quần áo, grayscale 28$\times$28
\cite{fashion}. Mỗi lớp giữ 600 ảnh validation bằng seed42, tạo train54k/val6k.
Mean/std chỉ tính từ train. Không augmentation để dễ quy chênh lệch về kiến trúc.
\FigureOrPlaceholder{assets/e1/dataset_overview.png}{Fashion-MNIST: ví dụ 10 lớp và phân bố train/validation}{Ảnh mẫu và số ảnh dùng trong thí nghiệm E1.}
\section{Kiến trúc và loss}
Linear flatten ảnh thành vector784, head10 tạo logits. MLP dùng 784--256--128--10,
ReLU/dropout0.2. CNN-A gồm conv32/64, pooling, FC128. CNN-B gồm conv32/64/128,
BN/ReLU/pool, adaptive average pool và dropout0.3. Cross-entropy áp dụng logits:
\begin{equation}\mathcal L=-\frac1N\sum_i\log\frac{\exp z_{i,y_i}}{\sum_c\exp z_{i,c}}.\end{equation}
Không softmax trước cross-entropy. Số tham số và thời gian đo từ mỗi model/run.
\section{Training protocol}
AdamW, LR $10^{-3}$, weight decay $10^{-4}$, batch128 (422 batch/epoch), tối đa 100 epoch, patience 4.
Early stopping khi validation loss không giảm ít nhất $10^{-4}$ trong 10 epoch liên tiếp.
Chọn checkpoint có validation loss thấp nhất của lần train hiện tại. CUDA AMP;
log xác nhận dữ liệu, parameters, logits và gradient trên GPU. Mỗi lần Run All train mới. Không dùng test
để chọn architecture hoặc early stop. Engine tự viết xử lý batch cuối và checkpoint.
\section{Kết quả và phân tích chờ hoàn thiện}
\ResultTable{e1}{So sánh bốn model E1; bảng full được xuất tự động.}
\FigureOrPlaceholder{assets/e1/comparison.png}{So sánh accuracy E1}{Accuracy validation của các full run E1.}
\FigureOrPlaceholder{assets/e1/learning_curves.png}{Learning curves của bốn model E1}{Loss và accuracy của Softmax, MLP, CNN-A và CNN-B; nét đứt train, nét liền validation.}
\FigureOrPlaceholder{assets/e1/confusion_test_pair_1.png}{Softmax và MLP: confusion matrix test}{Số ảnh và tỷ lệ theo hàng; mỗi lớp test có 1.000 ảnh.}
\FigureOrPlaceholder{assets/e1/confusion_test_pair_2.png}{CNN-A và CNN-B: confusion matrix test}{Cùng thang màu để so sánh bốn model.}
\FigureOrPlaceholder{assets/e1/failure_examples_test.png}{Ảnh lỗi test theo hàng model và cột lớp}{Một lỗi đầu tiên mỗi lớp/model, không đại diện toàn bộ lỗi.}
\Pending{Diễn giải metric test, confusion matrix và phân tích lỗi sau full run. So sánh
inductive bias, capacity và train/validation gap; kết luận phải dựa trên metric của lần chạy đầy đủ.}
Các cặp Shirt/T-shirt/Coat cần xem ảnh thật, kiểm tra texture/outline và độ mơ hồ lớp.
Phân tích tối thiểu một lợi ích và một giới hạn của từng family dựa trên evidence.

\IfFileExists{assets/e1/test_results.tex}{\begin{table}[H]\centering\caption{Metric test sau khi khóa protocol.}\resizebox{\linewidth}{!}{\begin{tabular}{lll}
\toprule
model & test\_accuracy & test\_macro\_f1 \\
\midrule
Softmax regression & 0.8427 & 0.8418 \\
MLP & 0.8945 & 0.8948 \\
CNN-A & 0.9166 & 0.9169 \\
CNN-B & 0.9197 & 0.9193 \\
\bottomrule
\end{tabular}}\end{table}}{}
